\documentclass{article}
\usepackage{graphicx} % Required for inserting images
%\usepackage[a4paper, total={6in, 8in}]{geometry}
\usepackage{amsmath, amssymb, empheq, mathtools, bbold, enumitem, hyperref, physics}
\usepackage{amsthm}
\usepackage[framemethod=TikZ]{mdframed}
%\usepackage{framed}
%\usepackage[amsthm,framed,thmmarks]{ntheorem}

\renewcommand{\bf}[1]{\textbf{#1}}
\renewcommand{\v}[1]{\boldsymbol{#1}}

%\newframedtheorem{frm-thm}{Theorem}
\theoremstyle{definition}
\newtheorem{theorem}{Theorem}
\newtheorem*{fact}{Fact}
\newtheorem*{note}{Note}
\newtheorem*{remark}{Remark}
\newtheorem{definition}{Definition}

\surroundwithmdframed[]{theorem}
\newenvironment{frmthm}{
	\begin{mdframed}
		\begin{theorem}
		}{		
		\end{theorem}
	\end{mdframed}
}

\newenvironment{factsummary}{
	{\large\textbf{Fact Summary}}\\
}{
}
\surroundwithmdframed{factsummary}

\newcommand{\xlrightarrow}{\xrightarrow{\hspace{0.8cm}}}
\newcommand{\reals}{\mathbb{R}}

\newcommand{\matadd}{+_{\text{mat}}}
\newcommand{\matdot}{\cdot_{\text{mat}}}
\newcommand{\vp}{\oplus}
\newcommand{\vm}{\odot}

\newcommand{\mxn}{m \times n}
\newcommand{\letmatrix}[2]{\mathcal{M}_{#1 \times #2}(\reals)}
\newcommand{\rowop}{\mathcal{R}}
%\newcommand{\span}[1]{\text{span(}#1\text{)}}


\newcommand{\bmat}[1]{\begin{bmatrix}#1\end{bmatrix}}
\newcommand{\detmat}[1]{
	\left|\hspace{0.5em}
	\begin{matrix}
		#1
	\end{matrix}
	\hspace{0.5em}\right|
}

% Ch. 4 Linear Transformations

\begin{document}
	\section{Linear Transformations}
	For $V, W$ which are real vector spaces, $L:V\rightarrow W$ is a \emph{linear transformation} if
	\begin{enumerate}[font=\sffamily\bfseries]
		\item $\forall v_1, v_2 \in V\\ L(v_1 + v_2) = L(v_1) + L(v_2)$
		\item $\forall v \in V, \forall c \in \reals\\ L(c\times v) = c \times L(v)$
	\end{enumerate}
	\subsection{Facts:}
	Let $L:V \rightarrow W$ be linear.
	\begin{enumerate}[font=\sffamily\bfseries]
		\item $L(0_V) = 0_W$
		\item $L(-v) = -L(v), \forall v \in V$
		\item For every vector subspace $V_0 of V$, $L(V_0)$ is a vector subspace of $W$.
		\begin{proof}
			Recall $L(V) = \{L(v) | v \in V\} \subseteq W$.\\
			Since $L(0_V) = 0_W, L(V) \neq \emptyset$. Therefore, it suffices to show that $L(V)$ is closed under vector addition and under scalar multiplication.\\\\
			Let $w_1, w_2 \in L(V)$.\\
			So, $w_1 = L(v_1)$ and $w_2 = L(v_2)$ for some $v_1, v_2 \in V$.
			Now $w_1 + w_2 = L(v_1) + L(v_2) = L(v_1 + v_2)$.
			Since $v_1 + v_2 \in V$, then $w_1 + w_2 \in L(V)$.\\
			For $c \in \reals$,\\
			$c \times w_1 = c \times L(v_1) = L(c \times v_1).\\
			\therefore c \times w_1 \in L(V).$
		\end{proof} 
	\end{enumerate}
	\subsection{Misc. Theorems}
	\begin{theorem}
		Let $V$ be a f-dim v.space and $W$ a v.subspace of $V$. Then $\dim(W) \leq \dim(V)$ and $\dim(W) = \dim(V) \iff W = V$.
		\begin{proof}
			Let $\phi = \{w_1, w_2, ..., w_d\}$ be a basis for $W$. So $\phi$ is a \emph{linearly independent} subset of $V$.\\
			Then there exists a basis of $V$ of the form $\{w_1,...,w_d; u_d,...,u_n\}$, when $n \geq d$.\\
			Since $\dim(W) = d$ and $\dim{V} = n$, then $d \leq n \implies \dim(W) \leq \dim(V)$.\\
			Furthermore $\dim(W) = \dim(V) \implies \phi$ is in fact a basis for $V$. In this case $W = V$.
		\end{proof}
	\end{theorem}
	\begin{theorem}
		Let $L, T:V \rightarrow W$ lin trans. Then the function $L+T:V \rightarrow W$ is also a linear transformation.
		\begin{proof} Show that $L+T$ satisfies both properties of a linear transformation.
			\begin{align*}
				(1) \quad (L+T)(v_1 + v_2) &= L(v_1 + v_2) + T(v_1 + v_2)\\
				&= L(v_1) + L(v_2) + T(v_1) + T(v_2)\\
				&= L(v_1) + T(v_1) + L(v_2) + T(v_2)\\
				&= (L+T)(v_1) + (L+T)(v_2)
			\end{align*}
			\begin{align*}
				(2) \quad c(L+T)(v) &= c(L(v) + T(v))\\
				&= c\cdot L(v) + c\cdot T(v)\\
				&= L(c \cdot v) + T(c \cdot v)\\
				&= (L+T)(c \cdot v)
			\end{align*}
		\end{proof}
	\end{theorem}
\end{document}